\aufzaehlungzwei {Als
\definitionsverweis {Einheitsnormalenfeld}{}{}
können wir vom Gradienten der beschreibenden Funktion, also von

\mavergleichskettedisp
{\vergleichskette
{ G(x,y)
}
{ =} { \begin{pmatrix}  2 \alpha x  \\ 2 \beta y    \end{pmatrix}
}
{ } {
}
{ } {
}
{ } {
}
}
{}{}{}
ausgehen, durch Normalisierung erhalten wir

\mavergleichskettedisp
{\vergleichskette
{ N(x,y)
}
{ =} {  { \frac{   1 }{  \sqrt{ \alpha^2 x^2 + \beta^2 y^2 }  } } \begin{pmatrix}   \alpha x  \\  \beta y    \end{pmatrix}
}
{ } {
}
{ } {
}
{ } {
}
}
{}{}{.}
Die Gauß-Abbildung ist also
\maabbeledisp {\varphi} {Y = \left\{  \begin{pmatrix}  x  \\ y   \end{pmatrix}   \in \R^2  \mid   \alpha x^2+\beta y^2 = 1        \right\}  } { S^1
} { \begin{pmatrix}  x  \\ y   \end{pmatrix} } { { \frac{   1 }{  \sqrt{ \alpha^2 x^2 + \beta^2 y^2 }  } } \begin{pmatrix}   \alpha x  \\  \beta y    \end{pmatrix}
} {.}
} {Wir behaupten, dass
\maabbeledisp {\psi} {S^1} {Y
} { \begin{pmatrix}  u  \\v   \end{pmatrix} } { { \frac{   1 }{  \sqrt{ { \frac{  u^2 }{ \alpha } }   + { \frac{  v^2 }{ \beta } } }  } }       \begin{pmatrix}   { \frac{   u  }{ \alpha } }   \\ { \frac{   v  }{ \beta } }    \end{pmatrix}
} {,}
eine Umkehrfunktion ist. Wegen

\mavergleichskettedisp
{\vergleichskette
{ { \frac{   1 }{   { \frac{  u^2 }{ \alpha } }   + { \frac{  v^2 }{ \beta } }  } } \left(  \alpha { \frac{  u^2 }{ \alpha^2 } } + \beta{ \frac{  v^2 }{ \beta^2 } }  \right)
}
{ =} { 1
}
{ } {
}
{ } {
}
{ } {
}
}
{}{}{}
liegt das Bild von $\psi$ in $Y$. Wegen

\mavergleichskettealign
{\vergleichskettealign
{  \varphi \left(  \psi \begin{pmatrix}  u  \\v   \end{pmatrix}  \right)
}
{ =} { \varphi \left(  { \frac{   1 }{  \sqrt{ { \frac{  u^2 }{ \alpha } }   + { \frac{  v^2 }{ \beta } } }  } }       \begin{pmatrix}   { \frac{   u  }{ \alpha } }   \\ { \frac{   v  }{ \beta } }    \end{pmatrix}    \right)
}
{ =} { \varphi \begin{pmatrix}   { \frac{   1 }{  \alpha    \sqrt{ { \frac{  u^2 }{ \alpha } }   + { \frac{  v^2 }{ \beta } } }  } }   u    \\ { \frac{   1 }{ \beta  \sqrt{ { \frac{  u^2 }{ \alpha } }  + { \frac{  v^2 }{ \beta } } }  } }  v       \end{pmatrix}
}
{ =} { { \frac{   1 }{  \sqrt{  \alpha^2 { \frac{  u^2 }{  \alpha^2   \left(   { \frac{  u^2 }{ \alpha } }  +  { \frac{  v^2 }{ \beta } }   \right)  } }  + \beta^2 { \frac{  v^2 }{  \beta^2 \left(   { \frac{  u^2 }{ \alpha } }  +  { \frac{  v^2 }{ \beta } }   \right)  } }       }  } }  \begin{pmatrix}  { \frac{   1 }{  \sqrt{ { \frac{  u^2 }{ \alpha } }   + { \frac{  v^2 }{ \beta } } }  } }   u    \\   { \frac{   1 }{   \sqrt{ { \frac{  u^2 }{ \alpha } } + { \frac{  v^2 }{ \beta } } }  } } v    \end{pmatrix}
}
{ =} { { \frac{   1 }{  \sqrt{ { \frac{  u^2 }{  \left(   { \frac{  u^2 }{ \alpha } }  +  { \frac{  v^2 }{ \beta } }   \right)  } }   + { \frac{  v^2 }{  \left(   { \frac{  u^2 }{ \alpha } }  +  { \frac{  v^2 }{ \beta } }   \right)  } }       }  } }  \begin{pmatrix}   { \frac{   1 }{  \sqrt{ { \frac{  u^2 }{ \alpha } }   + { \frac{  v^2 }{ \beta } } }  } }   u    \\   { \frac{   1 }{   \sqrt{ { \frac{  u^2 }{ \alpha } }   + { \frac{  v^2 }{ \beta } } }  } } v    \end{pmatrix}
}
}
{
\vergleichskettefortsetzungalign
{ =} { { \frac{    \sqrt{ { \frac{  u^2 }{ \alpha } }  +  { \frac{  v^2 }{ \beta } } } }{  \sqrt{    u^2+ v^2       }  } }  \begin{pmatrix}   { \frac{   1 }{  \sqrt{ { \frac{  u^2 }{ \alpha } }   + { \frac{  v^2 }{ \beta } } }  } }   u    \\   { \frac{   1 }{   \sqrt{ { \frac{  u^2 }{ \alpha } }   + { \frac{  v^2 }{ \beta } } }  } } v    \end{pmatrix}
}
{ =} { \begin{pmatrix}  u  \\v   \end{pmatrix}
}
{ } {
}
{ } {}
}
{}{}
und

\mavergleichskettealign
{\vergleichskettealign
{  \psi \left(  \varphi \begin{pmatrix}  x  \\ y   \end{pmatrix}  \right)
}
{ =} { \psi \begin{pmatrix}  { \frac{    \alpha x }{  \sqrt{\alpha^2 x^2+ \beta^2 y^2 }  } }   \\ { \frac{   \beta y }{  \sqrt{\alpha^2 x^2+ \beta^2 y^2  } }}    \end{pmatrix}
}
{ =} { { \frac{   1 }{  \sqrt{ { \frac{   \alpha x^2 }{  \alpha^2x^2 + \beta^2y^2  } }  + { \frac{   \beta y^2 }{  \alpha^2x^2 + \beta^2y^2  } }  }  } }  \begin{pmatrix}  { \frac{    x }{  \sqrt{\alpha^2 x^2+ \beta^2 y^2 }  } }   \\ { \frac{   y }{  \sqrt{\alpha^2 x^2+ \beta^2 y^2  } }}    \end{pmatrix}
}
{ =} { { \frac{   \sqrt{\alpha^2 x^2+ \beta^2 y^2 } }{  \sqrt{    \alpha x^2 +  \beta y^2  }  } }  \begin{pmatrix}  { \frac{    x }{  \sqrt{\alpha^2 x^2+ \beta^2 y^2 }  } }   \\ { \frac{   y }{  \sqrt{\alpha^2 x^2+ \beta^2 y^2  } }}    \end{pmatrix}
}
{ =} { \begin{pmatrix}  x  \\ y   \end{pmatrix}
}
}
{}
{}{}
sind die Abbildungen invers zueinander.
}

 [[Kategorie:Latexseite]]
